\documentclass[11pt]{article}
\usepackage[T1]{fontenc}
\usepackage{inter}
\renewcommand*\familydefault{\sfdefault}
\usepackage[dvipsnames]{xcolor}
\definecolor{gr}{RGB}{0,32,91}
\newcommand\ruleafter[1]{#1~\xrfill[.5ex]{1pt}[gr]}
\usepackage{geometry}
\geometry{
    a4paper,
    top=1.8cm,
    bottom=1in,
    left=2.5cm,
    right=2.5cm
}
\setcounter{secnumdepth}{0}
\pdfgentounicode=1
\usepackage{enumitem}
\setlist[itemize]{
    noitemsep,
    left=2em
}
\setlist[description]{itemsep=0pt}
\setlist[enumerate]{align=left}
\colorlet{icnclr}{gray}
\usepackage{titlesec}
\titlespacing{\subsection}{0pt}{*0}{*0}
\titlespacing{\subsubsection}{0pt}{*0}{*0}
\titleformat{\section}{\color{gr}\large\fontseries{black}\selectfont\uppercase}{}{}{\ruleafter}[\global\RemVStrue]
\titleformat{\subsection}{\fontseries{bold}\selectfont}{}{}{\rvs}
\titleformat{\subsubsection}{\color{gray}\fontseries{bold}\selectfont}{}{}{}
\usepackage{xhfill}
\newif\ifRemVS
\newcommand{\rvs}{
    \ifRemVS
    \vspace{-1.5ex}
    \fi
    \global\RemVSfalse
}
\usepackage{fontawesome5}
\usepackage[bookmarks=false]{hyperref}
\hypersetup{
    colorlinks=true,
    urlcolor=Sepia,
    pdftitle={Ho-min Park - CV},
}
\usepackage[page]{totalcount}
\usepackage{fancyhdr}
\pagestyle{fancy}
\renewcommand{\headrulewidth}{0pt}
\fancyhf{}
\usepackage{amsmath}
\usepackage{amsfonts}
\begin{document}
\definecolor{gr}{RGB}{0,32,91}
\titleformat{\section}{\color{gr}\large\fontseries{black}\selectfont\uppercase}{}{}{\ruleafter}[\global\RemVStrue]
%== HEADER ==%
\begin{center}
    {\fontsize{28}{28}\selectfont Ho-min Park} \\
    \bigskip
    {\color{icnclr}\faMapMarker} Houston, TX, USA \quad
    {\color{icnclr}\faEnvelope[regular]} \href{mailto:Ho-min.Park@bcm.edu}{Ho-min.Park@bcm.edu}, \href{mailto:powersimmani@gmail.com}{powersimmani@gmail.com} \quad{\color{icnclr}\faIcon{mobile-alt}} \href{tel:+82-10-4731+3162}{+82-10-4731-3162}
    \\
    {\color{icnclr}\faGithub} \href{https://github.com/powersimmani}{github.com/powersimmani}
    \quad{\color{icnclr}\faOrcid} \href{https://orcid.org/0000-0001-9937-8617}{0000-0001-9937-8617}
    \quad{\color{icnclr}\faLinkedin} \href{https://www.linkedin.com/in/park-ho-min-b46658a6}{LinkedIn}
    \quad{\color{icnclr}\faGlobe} \href{https://powersimmani.github.io}{powersimmani.github.io}
\end{center}

\section{About me}
Assistant Professor in Computer Science Engineering specializing in \textbf{biomedical and biological AI}, with extensive experience developing deep learning solutions for medical imaging (laparoscopy, MRI, microscopy), structural biology (CRISPR-Cas systems and anti-CRISPR proteins), and genomics. Demonstrated expertise in building interpretable, clinically oriented ML pipelines that have been validated in collaboration with surgeons, biologists, and environmental scientists. Author of peer-reviewed publications in high-impact venues including the \textit{Journal of Clinical Oncology}, \textit{International Journal of Surgery}, \textit{IEEE Transactions on Affective Computing}, and \textit{Transactions on Machine Learning Research}, holder of 2 patents, and contributor to multiple open-source bioinformatics tools. Currently an Assistant Professor at the Texas Children's Hospital Data Center / Baylor College of Medicine (Houston, TX), where his research focuses on advancing biomedical AI for pediatric and oncological applications, following a postdoctoral appointment at Ghent University Global Campus.

\bigskip
\noindent\textbf{Citation metrics (Google Scholar, May 2026):} 295 total citations; h-index 10; i10-index 10.

\section{Professional Experience}
\subsection{Assistant Professor\hfill \normalfont 2026 - Present}
\subsubsection{Texas Children's Hospital Data Center / Baylor College of Medicine, Houston, TX, USA}
\begin{itemize}
    \item[\checkmark] Leading research on machine learning for pediatric and oncological applications.
    \item[\checkmark] Building interpretable, clinically oriented deep learning pipelines in collaboration with clinicians and biomedical researchers.
\end{itemize}

\subsection{Postdoctoral Researcher\hfill \normalfont November 2025 - 2026}
\subsubsection{Ghent University Global Campus (GUGC), Incheon, South Korea}
\begin{itemize}
    \item[\checkmark] Conducted research on biomedical and biological AI under the supervision of Prof. Joris Vankerschaver.
    \item[\checkmark] Contributed to multi-institutional collaborations in surgical AI (peritoneal metastasis detection, laparoscopic video analysis) and genomics (data augmentation for genomic deep learning).
\end{itemize}

\section{Skills}
\subsection{Computational}
\subsubsection{Python \& Deep Learning Frameworks}
\begin{itemize}
    \item[\checkmark] Primary programming language, with expertise in AI model development using PyTorch and PyTorch Lightning.
    \item[\checkmark] Extensive experience implementing deep learning models for medical imaging, genomics, and multimodal data, focusing on MLOps tasks and integrating with Weights and Biases (WandB).
\end{itemize}
\subsubsection{GitHub}
\begin{itemize}
    \item[\checkmark] Extensively used for version control and project sharing, with experience in collaborative development and documentation.
    \item[\checkmark] Actively maintained multiple public repositories for AI research projects and educational materials.
\end{itemize}
\subsubsection{Weights and Biases (WandB)}
\begin{itemize}
    \item[\checkmark] Actively used for experiment tracking and model performance monitoring in MLOps tasks.
\end{itemize}
\subsection{Languages}
\subsubsection{Korean}
\begin{itemize}
    \item[\checkmark] Native speaker.
\end{itemize}
\subsubsection{English}
\begin{itemize}
    \item[\checkmark] \textit{Writing:} Demonstrated through technical documentation and research publications.
    \item[\checkmark] \textit{Reading:} Evident through the comprehension of technical and academic literature.
    \item[\checkmark] \textit{Speaking and Listening:} Applied in collaborative environments with international teams.
\end{itemize}

\section{Funded Projects}
\subsection{Depression Diagnosis \& Drug Adherence Research (2022 - 2025)}
\begin{itemize}
    \item[\checkmark] Developed advanced models for the MuSe-Stress and MuSe-Personalization challenges.
    \item[\checkmark] Implemented Transformer Encoders and novel feature engineering techniques.
    \item[\checkmark] Achieved 2nd place in MuSe-Personalisation 2023 and 3rd place in MuSe-Stress 2022.
    \item[\checkmark] Contributed to the field of affective computing with interpretable AI models.
    \item[\checkmark] \textbf{Funding:} National Research Foundation of Korea (Grant No.: NRF-2020K1A3A1A68093469), Department of Biotechnology in India (Grant No.: DBT/IC-12031(22)-ICD-DBT), and Ghent University Global Campus
\end{itemize}

\subsection{Environmental Monitoring Project (2020 - 2023)}
\begin{itemize}
    \item[\checkmark] Led development of MP-Net, an AI system for automated microplastics detection.
    \item[\checkmark] Designed and implemented a microplastics annotation tool for dataset creation.
    \item[\checkmark] Managed collaboration between environmental scientists and AI engineers.
    \item[\checkmark] \textbf{Funding:} Special Research Fund of Ghent University (Grant No.: 01N01718) and Ghent University Global Campus
\end{itemize}

\subsection{Smart Packaging System Development (2018 - 2021)}
\begin{itemize}
    \item[\checkmark] Developed and patented two AI-powered measurement systems for industrial conveyor belts.
    \item[\checkmark] Led implementation of 3D deep learning-based item classification system.
    \item[\checkmark] Designed integrated size and type monitoring system.
    \item[\checkmark] Collaborated with industry partners to deploy solutions in real factory environments.
    \item[\checkmark] \textbf{Funding:} Ministry of Trade, Industry, and Energy (MOTIE), Korea (Project No.: 10077285)
\end{itemize}

\section{Patents}
\begin{itemize}
    \item[\checkmark] Park, H., \& De Neve, W. (2020) An apparatus and method for measuring the size of high-speed objects and classifying their type on conveyors using light curtain sensors. Korean Intellectual Property Office. Application no.: 10-2020-0045098, Registration no.: 10-2273758.
    \item[\checkmark] Park, H., \& De Neve, W. (2018) Apparatus and method for measuring the size of high-speed boxes on conveyors using an RGB-D camera. Korean Intellectual Property Office. Application no.: 10-2018-0148118, Registration no.: 10-2066862.
\end{itemize}

\section{Publications}
\noindent\textit{Citation counts in italics are from Google Scholar (May 2026). First-author publications are marked with first-author position. Full publication list and metrics available at \href{https://scholar.google.com/citations?user=mFzW6LEAAAAJ&hl=ko}{Google Scholar}.}

\subsection{Journal Articles}
\begin{itemize}
    \item[\checkmark] Tozzi, F., \textbf{Park, H. M.}, Mousavi, S. A., Van Liefferinge, M., Moon, D., Fadaei, S., et al. (2026) Multimodal machine learning for staging laparoscopy: a combined image analysis and morphologic tool for the discrimination of peritoneal metastasis. \textit{International Journal of Surgery}, \textbf{112}(1): 373-383.
    \item[\checkmark] Tozzi, F., Mousavi, S. A., Van Liefferinge, M., Moon, D., \textbf{Park, H.}, Fadaei, S., et al. (2024) Development of a video-based deep learning model for differentiation of malignant and benign lesions during staging laparoscopy: Is the machine better than the expert? \textit{Journal of Clinical Oncology}, \textbf{42}(16\_suppl): e13616. \textit{(1 citation)}
    \item[\checkmark] Tozzi, F., \textbf{Park, H. M.}, Moon, D., Mousavi, S. A., Van Liefferinge, M., Fadaei, S., et al. (2024) Machine learning models for classification and morphological assessment of peritoneal lesions during staging laparoscopy. \textit{European Journal of Surgical Oncology}, \textbf{50}.
    \item[\checkmark] Lee, H., Ozbulak, U., \textbf{Park, H.}, Depuydt, S., De Neve, W., \& Vankerschaver, J. (2024) Assessing the reliability of point mutation as data augmentation for deep learning with genomic data. \textit{BMC Bioinformatics}, \textbf{25}(1): 170. \textit{(11 citations)}
    \item[\checkmark] \textbf{Park, H.} (first author), Kim, G., Oh, J., Van Messem, A., \& De Neve, W. (2024) Interpreting stress detection models using SHAP and attention for MuSe-Stress 2022. \textit{IEEE Transactions on Affective Computing}, \textbf{15}(01): 1-17. DOI: \href{https://doi.org/10.1109/TAFFC.2024.3488112}{10.1109/TAFFC.2024.3488112}. \textit{(9 citations)}
    \item[\checkmark] Ozbulak, U., Lee, H. J., Boga, B., Anzaku, E. T., \textbf{Park, H.}, Van Messem, A., et al. (2023) Know your self-supervised learning: A survey on image-based generative and discriminative training. \textit{Transactions on Machine Learning Research (TMLR)}. \textit{(71 citations)}
    \item[\checkmark] \textbf{Park, H.} (first author), Won, J., Park, Y., Anzaku, E. T., Vankerschaver, J., Van Messem, A., De Neve, W., \& Shim, H. (2023) CRISPR-Cas-Docker: Web-based in silico docking and machine learning-based classification of crRNAs with Cas proteins. \textit{BMC Bioinformatics}, \textbf{24}(1): 167. Springer. \textit{(5 citations)}
    \item[\checkmark] \textbf{Park, H.} (first author), Park, Y., Berani, U., Bang, E., Vankerschaver, J., Van Messem, A., De Neve, W., \& Shim, H. (2022) In silico optimization of RNA-protein interactions for CRISPR-Cas13-based antimicrobials. \textit{Biology Direct}, \textbf{17}(1). DOI: \href{https://doi.org/10.1186/s13062-022-00339-5}{10.1186/s13062-022-00339-5}. \textit{(13 citations)}
    \item[\checkmark] \textbf{Park, H. M.} (first author), Park, Y., Vankerschaver, J., Van Messem, A., De Neve, W., \& Shim, H. (2022) Rethinking protein drug design with highly accurate structure prediction of anti-CRISPR proteins. \textit{Pharmaceuticals}, \textbf{15}(3): 310. DOI: \href{https://doi.org/10.3390/ph15030310}{10.3390/ph15030310}. \textit{(25 citations)}
    \item[\checkmark] \textbf{Park, H.} (first author), Park, S., de Guzman, M. K., Baek, J. Y., Cirkovic Velickovic, T., Van Messem, A., \& De Neve, W. (2022) MP-Net: Deep learning-based segmentation for fluorescence microscopy images of microplastics isolated from clams. \textit{PLoS ONE}, \textbf{17}(6): e0269449. DOI: \href{https://doi.org/10.1371/journal.pone.0269449}{10.1371/journal.pone.0269449}. \textit{(32 citations)}
\end{itemize}

\subsection{Conference Articles}
\begin{itemize}
    \item[\checkmark] Mousavi, S. A., Tozzi, F., \textbf{Park, H.}, Anzaku, E. T., Van Liefferinge, M., Rashidian, N., et al. (2024) A Reference-Based Approach for Tumor Size Estimation in Monocular Laparoscopic Videos. \textit{International Workshop on Computational Mathematics Modeling in Cancer}. Springer. \textit{(2 citations)}
    \item[\checkmark] Singh, A. K., Kim, G., Kim, J., \textbf{Park, H.}, Choi, B. J., \& De Neve, W. (2025) RAMM: A Residual Attention Multimodal Model for Humor Detection. \textit{International Conference on Intelligent Human Computer Interaction}, 229-240. \textit{(3 citations)}
    \item[\checkmark] Nguyen, K. T., \textbf{Park, H.}, Oh, G., Vankerschaver, J., \& De Neve, W. (2025) Towards Improved Cervical Cancer Screening: Vision Transformer-Based Classification and Interpretability. \textit{IEEE 22nd International Symposium on Biomedical Imaging (ISBI)}, 1-5. \textit{(5 citations)}
    \item[\checkmark] \textbf{Park, H.} (first author), Yun, I., Kim, M., Nguyen, K. T., Van Messem, A., \& De Neve, W. (2024) Interpretable Rotator Cuff Tear Diagnosis Using MRI Slides with CAMscore and SHAP. \textit{Medical Imaging 2024: Computer-Aided Diagnosis}, Vol. 12927: 683-693. SPIE.
    \item[\checkmark] Singh, A. K., Kim, G., Kim, J., \textbf{Park, H.}, Choi, B. J., \& De Neve, W. (2024) Exploring Multimodal Approaches and Fusion Methods for CEO Social Attribute Prediction in 2024 MuSe-Perception. \textit{Proceedings of the 5th on Multimodal Sentiment Analysis Challenge and Workshop: Social Perception and Humor}, 36–44. Melbourne VIC, Australia: ACM. DOI: \href{https://doi.org/10.1145/3689062.3689081}{10.1145/3689062.3689081}. \textit{(2 citations)}
    \item[\checkmark] \textbf{Park, H. M.} (first author), Kim, G., Van Messem, A., \& De Neve, W. (2023) MuSe-Personalization 2023: Feature engineering, hyperparameter optimization, and transformer-encoder re-discovery. \textit{Proceedings of the 4th Multimodal Sentiment Analysis Challenge and Workshop: Mimicked Emotions, Humour and Personalisation}, 89–97. Ottawa ON, Canada: ACM. DOI: \href{https://doi.org/10.1145/3606039.3613104}{10.1145/3606039.3613104}. \textit{(5 citations)}
    \item[\checkmark] \textbf{Park, H.} (first author), Yun, I., Kumar, A., Singh, A. K., Choi, B. J., Singh, D., \& De Neve, W. (2022) Towards multimodal prediction of time-continuous emotion using pose feature engineering and a transformer encoder. \textit{Proceedings of the 3rd International on Multimodal Sentiment Analysis Workshop and Challenge}, 47–54. DOI: \href{https://doi.org/10.1145/3551876.3554807}{10.1145/3551876.3554807}. \textit{(13 citations)}
    \item[\checkmark] Kang, H., \textbf{Park, H.}, Ahn, Y., Van Messem, A., \& De Neve, W. (2021) Towards a quantitative analysis of class activation mapping for deep learning-based computer-aided diagnosis. \textit{Medical Imaging 2021: Image Perception, Observer Performance, and Technology Assessment}, Vol. 11599, SPIE. DOI: \href{https://doi.org/10.1117/12.2580819}{10.1117/12.2580819}. \textit{(18 citations)}
    \item[\checkmark] \textbf{Park, H.} (first author), Kang, B., Van Messem, A., \& De Neve, W. (2021) 3-D deep learning-based item classification for belt conveyors targeting packaging and logistics. \textit{Pattern Recognition, ICPR International Workshops and Challenges, Proceedings, Part VI}, 12666: 578–591. Cham: Springer. DOI: \href{https://doi.org/10.1007/978-3-030-68799-1_42}{10.1007/978-3-030-68799-1\_42}. \textit{(2 citations)}
    \item[\checkmark] Baek, J. Y., de Guzman, M. K., \textbf{Park, H.}, Park, S., Shin, B., Cirkovic Velickovic, T., Van Messem, A., \& De Neve, W. (2021) Developing a segmentation model for microscopic images of microplastics isolated from clams. \textit{Pattern Recognition, ICPR International Workshops and Challenges, Proceedings, Part VI}, 12666: 86–97. Cham: Springer. DOI: \href{https://doi.org/10.1007/978-3-030-68780-9_9}{10.1007/978-3-030-68780-9\_9}. \textit{(4 citations)}
    \item[\checkmark] Kim, M., \textbf{Park, H.}, Kim, J. Y., Kim, S. H., Hoeke, S., \& De Neve, W. (2020) MRI-based diagnosis of rotator cuff tears using deep learning and weighted linear combinations. \textit{Machine Learning for Healthcare Conference (MLHC)}, 292-308. \textit{(19 citations)}
    \item[\checkmark] \textbf{Park, H.} (first author), Van Messem, A., \& De Neve, W. (2020) Item measurement for logistics-oriented belt conveyor systems using a scenario-driven approach and automata-based control design. \textit{IEEE 7th International Conference on Industrial Engineering and Applications (ICIEA 2020)}, Bangkok and Phuket. \textit{(8 citations)}
    \item[\checkmark] \textbf{Park, H. M.} (first author), Van Messem, A., \& De Neve, W. (2019) Box-Scan: An efficient and effective algorithm for box dimension measurement in conveyor systems using a single RGB-D camera. \textit{The 7th IIAE International Conference on Industrial Application Engineering}, Kitakyushu, Japan. \textit{(16 citations)} \textbf{[Best Presentation Award]}
\end{itemize}

\section{Awards}
\begin{itemize}
    \item[\checkmark] \textbf{Second Place}, MuSe-Personalisation Challenge 2023 — international competition at ACM Multimedia, for contributions to multimodal sentiment analysis.
    \item[\checkmark] \textbf{Third Place}, MuSe-Stress Challenge 2022 — international competition at ACM Multimedia, for contributions to multimodal sentiment analysis.
    \item[\checkmark] \textbf{Best Presentation Award}, 7th IIAE International Conference on Industrial Application Engineering, 2019.
\end{itemize}

\section{Education}
\subsection{Ph.D. in Computer Science Engineering\hfill \normalfont 2018 - 2025}
\subsubsection{Ghent University, Belgium}
\begin{itemize}
    \item[\checkmark] Dissertation: ``Advancing Detection through Deep Learning for Industrial, Environmental, and Health Applications''
    \item[\checkmark] Advisors: Prof. Wesley De Neve, Prof. Arnout Van Messem
\end{itemize}

\subsection{M.S. in Computer Engineering\hfill \normalfont 2016 - 2018}
\subsubsection{Ajou University, Suwon, South Korea}
\begin{itemize}
    \item[\checkmark] Thesis: ``Cited count prediction modeling from large citation network''
\end{itemize}

\subsection{B.S. in Computer Science and Engineering\hfill \normalfont 2010 - 2016}
\subsubsection{Ajou University, Suwon, South Korea}
\begin{itemize}
    \item[\checkmark] Major in Computer Science and Engineering
\end{itemize}

\section{Mentoring and Educational Leadership}
\subsection{Founder and Lead Instructor, AI Vacation School (AIVS)}
\begin{itemize}
    \item[\checkmark] Founded and led the \textit{AI Vacation School (AIVS)} at the Center for Biosystems and Biotech Data Science, Ghent University Global Campus, as a volunteer initiative to provide intensive AI education to undergraduate students from 2021 to 2026, supporting over 30 participants across summer and winter cohorts. [\href{https://github.com/powersimmani/AIVS}{Program Repository}, \href{https://powersimmani.github.io/AIVS_Lecture_Slides/}{Lecture Slides}]
    \item[\checkmark] Mentored participants who have advanced to top graduate programs and research positions, including KAIST, Seoul National University, UNIST, Ghent University (Belgium), and Scripps Research Institute (San Diego, CA, USA).
    \item[\checkmark] Co-authored peer-reviewed publications with alumni in venues such as \textit{IEEE Transactions on Affective Computing}, \textit{BMC Bioinformatics}, \textit{PLoS ONE}, \textit{SPIE Medical Imaging}, and \textit{ACM Multimodal Sentiment Analysis Workshop}, demonstrating the program's effectiveness in producing research-active students.
    \item[\checkmark] Developed and released open-source educational and research artifacts in collaboration with participants, including \textit{CRISPR-Cas-Docker} (\href{http://www.crisprcasdocker.org/}{web service}), \textit{MP-Net} segmentation toolkit, \textit{Microplastics Annotation Package}, and the \textit{MPSet} dataset on Kaggle.
    \item[\checkmark] Presented the program at the \textit{IGC Research Showcase 2025, AI in Education} session, sharing pedagogical insights with the four universities of Incheon Global Campus (SUNY Korea, University of Utah Asia Campus, George Mason University Korea, Ghent University Global Campus). [\href{https://www.dropbox.com/scl/fi/3cdkfs2v846sj1g000jo6/IGC-Research-Showcases-Ho-min-Park.pdf?rlkey=s7qrhj4vva4u2rpboeccoa5aq&st=1mlzmm6j&dl=0}{Presentation}]
    \item[\checkmark] Provided extensive individualized mentoring, including portfolio development guidance, reference letter writing for graduate school and industry applications, and ongoing career advising for alumni.
\end{itemize}

\section{References}
\textcolor{Sepia}{References available upon request.}

\end{document}
